\section{Appendix: Liturgical Resolution}
This is the Twenty-seventh Sunday in Ordinary Time, Year A, 4 October 2026, for the dioceses of the United States under the General Roman Calendar and U.S. national proper: a Sunday in Ordinary Time, green, Lectionary 139. The governing Missal is the \work{Roman Missal, Third Edition}, U.S. English edition implemented in 2011, with the 2008 emended Latin third typical edition as base. The U.S. \work{Lectionary for Mass}, second typical edition (1998/2001), Volume I, controls the readings. The ordinary Sunday cycle is A; the independently resolved adjacent weekday cycle II does not supply this Sunday's lessons.\footnote{USCCB, \work{Liturgical Calendar for the Dioceses of the United States of America 2026}, printed pp. 5,40; dated readings for 4 October 2026, Lectionary 139.}

The national calendar assigns this Sunday rather than the occurring memorial of Francis of Assisi. No local dedication, titular solemnity, patronal observance or religious calendar has been resolved. The permanent formula is PC-S53-A, parent PC-S53. The common Missal evidence belongs to the edition's Week 27 owner; the Year A leaf owns this reading composition. Neither another cycle nor another Missal supplies an additional proper.

\begin{branchtable}
\bid{entrance-missal-or-permitted-chant} & GIRM 48; singing and no-singing routes & permitted & Entrance & Missal antiphon studied; no concrete chant selected. Its recitation is required when there is no entrance singing.\\
\bid{psalm-lectionary-or-permitted-sung-form} & GIRM 61 & permitted & Psalm & Appointed psalm and Isaiah response studied; sung alternatives unselected.\\
\bid{acclamation-sung-or-omitted} & GIRM 62--64 & conditional & Acclamation & Appointed adapted verse studied. GIRM 63(c)'s one-reading condition does not authorize its omission at this three-reading Sunday.\\
\bid{communion-lamentations-or-corinthians} & Week 27 Missal's explicit alternative & appointed alternative & Communion & Both studied; neither a local selection.\\
\bid{communion-missal-or-permitted-chant} & GIRM 87--88 & permitted & Communion & Chant classes remain open; no additional concrete chant selected.\\
\bid{preface-eucharistic-prayer} & GIRM 364--365 & permitted & Eucharistic Prayer & No unique Week 27 choice. Any selection must retain the applicable Preface coupling; Eucharistic Prayer IV has its own invariable Preface.\\
\end{branchtable}
Sprinkling in place of the Penitential Act (GIRM 51), a preparation chant (74), and any optional concluding blessing remain unselected. Gloria, Creed, homily and Universal Prayer belong to the Sunday's ordinary requirements (53,65--71); they do not create additional closed proper texts. No local musical enactment is asserted.

\section{Appendix: Scope and Qualifications}
This study develops two complementary readings of this formulary from its biblical contexts and the specified patristic and saintly loci. Their connections across the whole Mass are editorial synthesis. Isaiah's relation to the Gospel belongs to the Lectionary's correlated selection; Philippians follows its semi-continuous course, and the common prayers recur across cycles. No ancient writer is credited with an exposition of this entire modern formulary.

The English Scripture is public-domain Douay--Rheims Challoner study text, not the approved proclaimed U.S. English. Protected Lectionary and ICEL wording, full restricted Latin prayers, and newly composed prayer translations are not reproduced. The canonical Week 27 evidence rests on an identified 2002 digital Latin reproduction, official ICEL antiphon excerpt and approved FDLC prayer excerpts, with the limited evidence of the 2008 variation list in \work{Notitiae} 44, pp. 367--372. A complete U.S. altar-book and 2008 page collation remain absent; the exact offerings English remains uncollated. The FDLC Collect's older conclusion is not adopted.

Reception rests on the specified historical translations and digital witnesses. Augustine's psalm exposition and O'Sullivan's Bellarmine are abridged; neither licenses claims about omitted Latin passages. No fresh Greek or Latin critical collation is claimed. Direct, securely attributed saintly commentary on the exact Esther and Lamentations loci was not established in the bounded research. Their biblical contexts carry the exposition; an uncertainly attributed Lamentations commentary is not cited as secure Aquinas. Aquinas's \work{Summa} and Leo's sermon illuminate sacramental meaning, rather than interpreting the modern prayers as such.

The Catholic Encyclopedia chronology witnesses are historical reference judgments; the generated dates preserve relation types, competing proposals and unresolved states. The study does not claim universal reception coverage, ecclesiastical approval, exact liturgical enactment, or authority to replace the governing books. Source bindings, the interpretation audit, chronology record and shared-owner audit retain the detailed evidence and its limits. Research and author verification are distinct from the subsequent independent reviews of the rendered publications.

\section{References}
\begin{itemize}
\item \work{Holy Bible}, Douay--Rheims, Challoner revision; Project Gutenberg 1581 witness. Esther 4,13--14; Isaiah 5; Psalm 79; Lamentations 3; Matthew 21; John 15; 1 Corinthians 10; Philippians 4. Public-domain study English.
\item Augustine, \work{Tractates on John} 86--87, trans. John Gibb and James Innes, \work{Nicene and Post-Nicene Fathers}, first series, vol. 7 (NPNF1-7); and \work{Exposition of Psalm} 79, §§1--11, NPNF1-8, English heading Psalm 80. Historical English in CCEL witnesses; the psalm exposition is abridged.
\item John Chrysostom, \work{Homilies on Matthew} 68.1--2, NPNF1-10; \work{Homilies on First Corinthians} 24.1--5, especially 24.4, NPNF1-12; \work{Homilies on Philippians} 14, NPNF1-13. Historical English CCEL witnesses.
\item Robert Bellarmine, \work{A Commentary on the Book of Psalms}, trans. John O'Sullivan (Dublin, 1866), Psalm 79. Abridged historical English, retained e-Catholic 2000 witness.
\item Thomas Aquinas, \work{Summa theologiae} III, q. 79, a. 1, body and replies; trans. Fathers of the English Dominican Province, Project Gutenberg 19950 witness.
\item Leo the Great, Sermon 63, \work{On the Passion XII}, VII and note 1063, trans. Charles Lett Feltoe, NPNF2-12, CCEL witness.
\item \work{Ordo lectionum Missae}, second typical edition (1981), General Introduction 23,66--68,106--107; approved English hosted by the Liturgy Office of England and Wales.
\item USCCB, \work{Liturgical Calendar for the Dioceses of the United States of America 2026}, printed pp. 5,40; \href{https://bible.usccb.org/bible/readings/100426.cfm}{readings for 4 October 2026}, Lectionary 139. Official dated witnesses inspected 28 September 2026.
\item \work{General Instruction of the Roman Missal}, U.S. English, chapters II and VII, nos. 48,51,53--54,61--71,74,87--90,364--365; official USCCB states inspected 28 September 2026.
\item \work{Missale Romanum}, third typical edition (2002), secondary digital reproduction, physical pp. 290--291; ICEL \work{Antiphonary} (2010), official CBCEW excerpt, printed p. 80; FDLC, \work{Mystagogical Reflections}, Ordinary Time 24--30 (2016), pp. 12--13. Exact-text and rights limits are stated in the scope appendix.
\item \work{Notitiae} 44, nos. 503--504 (2008), pp. 367--372, especially p. 371: bounded variation notice, not proof of complete edition identity.
\item \work{Catholic Encyclopedia}: Arthur McMahon, ``Esther'' (1909); Charles Souvay, ``Isaias'' (1910); Ferdinand Prat, ``St. Paul'' (1911); Achille Vander Heeren, ``Epistle to the Philippians'' (1911); Cornelius Aherne, ``Epistles to the Corinthians'' (1908); Leopold Fonck, ``Gospel of St. John'' (1910); Anthony Maas, ``Chronology of the Life of Jesus Christ'' (1910); Alfred Durand, ``The Synoptics'' (1912); E. Jacquier, ``Gospel of St. Matthew'' (1911). Chronological discussions identified in the canonical chronology record.
\item NABRE introductions to Psalms and Lamentations, official USCCB witnesses: broad psalm composition limits and the setting, authorship and grouping of Lamentations. Used for chronology and orientation; protected Bible translation not reproduced.
\end{itemize}
